\documentclass[11pt,a4paper]{article}
\usepackage[utf8]{inputenc}
\usepackage[T1]{fontenc}
\usepackage[french]{babel}
\usepackage{amsmath,amssymb,amsfonts}
\usepackage{hyperref}
\usepackage{booktabs}
\usepackage{geometry}
\geometry{margin=2.5cm}

\title{\textbf{Yggdrasil Engine --- Formules} \\[0.5em]
\large Sources, modifications et mapping vers le mycelium}
\author{Sky \& Claude (Opus 4.6) --- Sessions 6--35, 23 f\'ev -- 7 avril 2026}
\date{}

\begin{document}
\maketitle

% ============================================================================
\section{Architecture des strates}
% ============================================================================

\begin{center}
\begin{tabular}{clll}
\toprule
\textbf{Strate} & \textbf{Nom} & \textbf{Contenu} & \textbf{Mycelium?} \\
\midrule
S6 & Ind\'ecidable & G\"odel, Halting problem & Non \\
S5 & Presque ind\'ecidable & & Non \\
S4 & Logique sup\'erieure & & Non \\
S3 & Conjectures & Riemann, P=NP & Non \\
S2 & R\'ecursion\textsuperscript{2} & & Non \\
S1 & Structures r\'ecursives & & Non \\
\midrule
\textbf{S0} & \textbf{Concepts prouv\'es} & \textbf{65,026 concepts OpenAlex} & \textbf{Oui} \\
S-1 & M\'etiers & physics, biology, eng. & Oui \\
S-2 & Glyphes & $=$, $+$, $\int$, $\Sigma$, $\partial$ & Oui \\
\bottomrule
\end{tabular}
\end{center}

Le mycelium (r\'eseau de co-occurrences) vit dans le \textbf{sol} (S-2 \`a S0).
Au-dessus de S0 = l'arbre (conjectures, abstractions). Pas de mycelium.

% ============================================================================
\section{Formules utilis\'ees --- Sources originales}
% ============================================================================

% ----------------------------------------------------------------------------
\subsection{Fitness $\eta_i$ --- Wang, Song \& Barab\'asi (2013)}
% ----------------------------------------------------------------------------

\textbf{Source:} Wang, D., Song, C., \& Barab\'asi, A.-L. (2013).
\textit{Quantifying Long-term Scientific Impact}.
Science, 342(6154), 127--132.
\href{https://doi.org/10.1126/science.1237825}{doi:10.1126/science.1237825}

\textbf{Formule originale:}
\begin{equation}
c_i(t) = \eta_i \cdot m \cdot \sum_j (c_j + 1) \cdot P_i(t - t_i)
\end{equation}
o\`u $P_i(t)$ est un vieillissement log-normal:
\begin{equation}
P_i(t) = \frac{1}{t} \exp\!\left(-\frac{(\ln t - \mu_i)^2}{2\sigma_i^2}\right)
\end{equation}

\textbf{Adaptation Yggdrasil} (\texttt{engine/core/scisci.py:15}):
\begin{equation}
\eta_i \approx \frac{\text{citations\_totales}}{\sum_{t=1}^{\text{age}} P_i(t)}
\quad\text{avec } \mu=1, \sigma=1
\end{equation}
Simplification: on estime $\eta_i$ comme ratio citations r\'eelles / citations attendues
selon l'\^age, sans le terme de pr\'ef\'erence cumulative.

% ----------------------------------------------------------------------------
\subsection{D-index (Disruption) --- Wu, Wang \& Evans (2019)}
% ----------------------------------------------------------------------------

\textbf{Source:} Wu, L., Wang, D., \& Evans, J. A. (2019).
\textit{Large Teams Develop and Small Teams Disrupt Science and Technology}.
Nature, 566, 378--382.
\href{https://doi.org/10.1038/s41586-019-0941-9}{doi:10.1038/s41586-019-0941-9}

Formule originale (Funk \& Owen-Smith, 2017):
\begin{equation}
D = \frac{n_i - n_j}{n_i + n_j + n_k}
\end{equation}
\begin{itemize}
\item $n_i$ = papers citant le focal MAIS PAS ses r\'ef\'erences
\item $n_j$ = papers citant le focal ET ses r\'ef\'erences
\item $n_k$ = papers citant SEULEMENT les r\'ef\'erences
\end{itemize}
$D \to +1$: disruptif (nouveau paradigme).
$D \to -1$: developmental (extension).

\textbf{Adaptation Yggdrasil} (\texttt{engine/core/scisci.py:55}):
Utilis\'e tel quel. Sert \`a d\'etecter les trous techniques (type A: $|D|$ bas)
et perceptuels (type C: $|D|$ haut mais citations basses).

% ----------------------------------------------------------------------------
\subsection{Z-score d'atypicit\'e --- Uzzi et al. (2013)}
% ----------------------------------------------------------------------------

\textbf{Source:} Uzzi, B., Mukherjee, S., Stringer, M., \& Jones, B. (2013).
\textit{Atypical Combinations and Scientific Impact}.
Science, 342(6157), 468--472.
\href{https://doi.org/10.1126/science.1240474}{doi:10.1126/science.1240474}

\begin{equation}
z = \frac{f_{\text{obs}} - \mu_{\text{random}}}{\sigma_{\text{random}}}
\end{equation}
$z \ll 0$: combinaison tr\`es atypique (novel).
$z \gg 0$: combinaison conventionnelle.

\textbf{Adaptation Yggdrasil} (\texttt{engine/core/scisci.py:83}):
Utilis\'e tel quel. Sert au score de trou conceptuel (type B):
plus le z-score est n\'egatif, plus la paire est un trou potentiel.

% ----------------------------------------------------------------------------
\subsection{Q-factor --- Sinatra et al. (2016)}
% ----------------------------------------------------------------------------

\textbf{Source:} Sinatra, R., Wang, D., Deville, P., Song, C., \& Barab\'asi, A.-L. (2016).
\textit{Quantifying the Evolution of Individual Scientific Impact}.
Science, 354(6312), aaf5239.
\href{https://doi.org/10.1126/science.aaf5239}{doi:10.1126/science.aaf5239}

\begin{equation}
\log(c_i) = \log(Q_i) + \log(p_i)
\end{equation}
$Q$ = qualit\'e constante d'un chercheur. $p$ = chance (random).

\textbf{Adaptation Yggdrasil} (\texttt{engine/core/scisci.py:107}):
\begin{equation}
Q \approx \text{median}\!\left(\log(c_i + 1)\right)
\end{equation}
Simplification: m\'ediane des log-citations comme proxy de qualit\'e.

% ============================================================================
\section{Formules de d\'etection de trous}
% ============================================================================

\textbf{Fichier:} \texttt{engine/core/holes.py}

% ----------------------------------------------------------------------------
\subsection{Type A --- Trou Technique}
% ----------------------------------------------------------------------------

Tout le monde SAIT o\`u aller, personne ne PEUT.
Exemples: Poincar\'e (Hamilton bloqu\'e 20 ans), Higgs (besoin du LHC).

\begin{equation}
\text{Score}_A = \text{Production}(t) \times (1 - |\Delta\eta / \Delta t|) \times (1 - |D|)
\end{equation}
Production \'elev\'ee $\times$ fitness stagnante $\times$ D-index bas (developmental).

% ----------------------------------------------------------------------------
\subsection{Type B --- Trou Conceptuel}
% ----------------------------------------------------------------------------

Personne n'a l'ID\'EE de connecter. Le vide est invisible.
Exemples: GANs (game theory $\times$ deep learning), CRISPR, AlphaFold.

\begin{equation}
\text{Score}_B = \text{Activity}(A) \times \text{Activity}(B) \times (1 - \text{CoOcc}(A,B)) \times \frac{|z|}{10}
\end{equation}

% ----------------------------------------------------------------------------
\subsection{Type C --- Trou Perceptuel}
% ----------------------------------------------------------------------------

L'outil EXISTE, personne n'y CROIT.
Exemples: mRNA (Karik\'o rejet\'ee 30 ans), H.\ pylori, quasicristaux.

\begin{equation}
\text{Score}_C = \eta_i \times |D| \times \left(1 - \frac{c_i(t)}{c_i^{\text{expected}}}\right)
\end{equation}

% ============================================================================
\section{Topologie du mycelium}
% ============================================================================

\textbf{Fichier:} \texttt{engine/pipeline/mycelium\_full.py} (24 briques, 456 tests)

% ----------------------------------------------------------------------------
\subsection{Meshedness $\alpha$ (brique 1)}
% ----------------------------------------------------------------------------

\textbf{Source:} Bebber, D. P., Hynes, J., Darrah, P. R., Boddy, L., \& Fricker, M. D. (2007).
\textit{Biological solutions to transport network design}.
Proc.\ R.\ Soc.\ B, 274, 2307--2315.
\href{https://doi.org/10.1098/rspb.2007.0459}{doi:10.1098/rspb.2007.0459}

\begin{equation}
\alpha = \frac{L - N + 1}{2N - 5}
\end{equation}
$L$ = ar\^etes, $N$ = n\oe uds. $\alpha = 0$: arbre pur. $\alpha = 1$: r\'eseau planaire maximal.

\textbf{Adaptation Yggdrasil:} utilis\'e tel quel. Le $\alpha$ mesure la r\'esilience
du mycelium dans le sol (S-2 \`a S0). Sert au delta $\alpha$ dans les bougies OHLC.

% ----------------------------------------------------------------------------
\subsection{Betweenness Centrality (brique 5)}
% ----------------------------------------------------------------------------

\textbf{Source:} Freeman, L. C. (1977).
\textit{A Set of Measures of Centrality Based on Betweenness}.
Sociometry, 40(1), 35--41.
\href{https://doi.org/10.2307/3033543}{doi:10.2307/3033543}

\begin{equation}
BC(v) = \sum_{s \neq v \neq t} \frac{\sigma_{st}(v)}{\sigma_{st}}
\end{equation}
$\sigma_{st}$ = chemins les plus courts entre $s$ et $t$.
$\sigma_{st}(v)$ = ceux passant par $v$.

\textbf{Adaptation Yggdrasil:} utilis\'e via NetworkX (normalis\'e). Un n\oe ud \`a BC \'elev\'e
est un pont (P1). Corr\`ele avec le flux r\'eel en mycelium (Oyarte Galvez 2025).

% ----------------------------------------------------------------------------
\subsection{Physarum (brique 10)}
% ----------------------------------------------------------------------------

\textbf{Source:} Tero, A., Takagi, S., Saigusa, T., et al. (2010).
\textit{Rules for biologically inspired adaptive network design}.
Science, 327(5964), 439--442.
\href{https://doi.org/10.1126/science.1177894}{doi:10.1126/science.1177894}

\begin{equation}
\frac{dD_{ij}}{dt} = f(|Q_{ij}|) - \gamma D_{ij}
\end{equation}
$D_{ij}$ = conductivit\'e de l'hyphe. $Q_{ij}$ = flux. $\gamma$ = d\'ecroissance.
$f(|Q|) = |Q|^\mu$ avec $\mu = 1$ (shortest path) ou $\mu < 1$ (loops tolerated).

\textbf{Adaptation Yggdrasil:} les hyphes vivantes/mortes du Physarum servent
au delta 6 (redistribution des flux) dans les bougies OHLC.

% ----------------------------------------------------------------------------
\subsection{Laplacien spectral (positionnement)}
% ----------------------------------------------------------------------------

\textbf{Source:} Chung, F. R. K. (1997).
\textit{Spectral Graph Theory}. AMS.

Laplacien normalis\'e:
\begin{equation}
\mathcal{L} = D^{-1/2} L D^{-1/2} = I - D^{-1/2} A D^{-1/2}
\end{equation}
$A$ = matrice de co-occurrence (log-transform\'ee). $D$ = matrice des degr\'es.

Vecteur de Fiedler (2\`eme plus petite valeur propre):
\begin{equation}
\mathcal{L} \mathbf{v}_2 = \lambda_2 \mathbf{v}_2
\end{equation}

\textbf{Adaptation Yggdrasil} (\texttt{engine/topology/cooccurrence\_scan.py:225}):
\begin{equation}
\text{position}_x = v_2[i], \quad \text{position}_z = v_3[i]
\end{equation}
Les 2\`eme et 3\`eme vecteurs propres donnent les coordonn\'ees spectrales.
Appliqu\'e sur la matrice de co-occurrence \textbf{log-transform\'ee}: $W = \log(1 + C)$
o\`u $C$ est la matrice de co-occurrence brute.

Positionne TOUT le sol: glyphes (S-2), m\'etiers (S-1), concepts (S0).

% ----------------------------------------------------------------------------
\subsection{5 curseurs Lehmann (identification d'esp\`ece, session 8)}
% ----------------------------------------------------------------------------

\textbf{Source:} Lehmann, S., Lautrup, B., \& Jackson, A. D. (2019).
\textit{Citation networks in high energy physics}.
Phys.\ Rev.\ E, 68, 026113.

5 curseurs mesur\'es sur le graphe de co-occurrences:
\begin{align}
BA &= \text{angle de bifurcation moyen} \\
IL &= \text{internode length (longueur entre branchements)} \\
D  &= \text{diam\`etre du r\'eseau local} \\
D_b &= \text{dimension fractale (Minkowski)} \\
L  &= \text{lacunarit\'e (h\'et\'erog\'en\'eit\'e spatiale)}
\end{align}

\textbf{Adaptation Yggdrasil} (\texttt{engine/topology/species\_identifier.py}):
9 esp\`eces spectrales (K=9 sur 65K):
MatSci/Chem, Geo/Env, Medicine, Psych/Business, CS/Math,
Bio/Botany, Humanities/PolSci, CellBio/Anatomy, Physics/Optics.

% ----------------------------------------------------------------------------
\subsection{Impact Scale (session 31)}
% ----------------------------------------------------------------------------

\textbf{Fichier:} \texttt{engine/topology/impact\_scale.py}

Score d'impact normalis\'e 0--10 pour chaque paper:
\begin{equation}
\text{impact}(p) = \sum_{c \in \text{concepts}(p)} \text{rarity}(c) \times \text{spread}(c)
\end{equation}
o\`u rarity $= 1 / \log(\text{works\_count})$ et spread $=$ nombre d'esp\`eces touch\'ees.

Classification par tiers: caillou (0--3), pierre (3--5), rocher (5--7), m\'et\'eorite (7--10).

% ----------------------------------------------------------------------------
\subsection{Blind test V2 (session 11)}
% ----------------------------------------------------------------------------

\textbf{M\'ethode:} cutoff 2015 sur les 65,026 concepts. Train sur donn\'ees $\leq$ 2015,
predict les trous P4 qui se ferment apr\`es 2015.

\textbf{R\'esultat:} Mann-Whitney $U$ test, $p = 3.4 \times 10^{-12}$, Cohen's $d = 0.44$.
\begin{equation}
U = n_1 n_2 + \frac{n_1(n_1 + 1)}{2} - R_1
\end{equation}

Le moteur d\'etecte des trous structurels qui se ferment significativement plus
souvent que le hasard ($p < 10^{-11}$).

% ============================================================================
\section{Formules Sedov-Taylor --- Impact m\'et\'eorite [INVALID\'E session 33]}
% ============================================================================

\textbf{INVALID\'E session 33 (2 avril 2026).}
R$^2$ holdout G\"odel = $-5.74$. Sedov mod\'elise une explosion dans un gaz 3D,
le myc\'elium est une mare 2D. Remplac\'e par le mod\`ele caillou-dans-la-mare
(sections 9--13).

% ----------------------------------------------------------------------------
\subsection{Formules originales (physique)}
% ----------------------------------------------------------------------------

\textbf{Sources:}
\begin{itemize}
\item Taylor, G. I. (1950). \textit{The Formation of a Blast Wave by a Very Intense Explosion}.
Proc.\ R.\ Soc.\ A, 201(1065), 159--174.
\item Sedov, L. I. (1946). \textit{Propagation of strong shock waves}.
J.\ Appl.\ Math.\ Mech., 10, 241--250.
\item von Neumann, J. (1947). \textit{The point source solution}. In: Blast Wave (report).
\end{itemize}

\textbf{Rayon du front de choc:}
\begin{equation}
\boxed{R(t) = \beta \left(\frac{E t^2}{\rho_0}\right)^{1/5}}
\end{equation}
\begin{itemize}
\item $E$ = \'energie totale inject\'ee (J)
\item $\rho_0$ = densit\'e du milieu ambiant (kg/m$^3$)
\item $t$ = temps depuis l'explosion
\item $\beta$ = constante sans dimension:
  $\beta \approx 1.033$ pour $\gamma = 7/5$ (air),
  $\beta \approx 1.152$ pour $\gamma = 5/3$ (gaz monoatomique)
\end{itemize}

\textbf{Vitesse du choc:}
\begin{equation}
\dot{R}(t) = \frac{2}{5} \beta \left(\frac{E}{\rho_0}\right)^{1/5} t^{-3/5}
\end{equation}
Le choc d\'ec\'el\`ere en $t^{-3/5}$.

\textbf{Conditions post-choc (Rankine-Hugoniot, \`a $\xi = 1$):}
\begin{align}
V(1) &= \frac{2}{\gamma + 1} \\
G(1) &= \frac{\gamma + 1}{\gamma - 1} \\
Z(1) &= \frac{2}{\gamma + 1}
\end{align}
Pour $\gamma = 7/5$: densit\'e post-choc $= 6\times$ la densit\'e ambiante.

\textbf{Conservation d'\'energie:}
\begin{equation}
E = 4\pi \rho_0 R^3 \dot{R}^2 \left[
2\int_0^1 G(\xi) V^2(\xi) \xi^2 \, d\xi
+ \frac{2}{\gamma - 1} \int_0^1 Z(\xi) \xi^2 \, d\xi
\right]
\end{equation}
Partition pour $\gamma = 7/5$: $\sim$30\% cin\'etique, $\sim$70\% thermique.

% ----------------------------------------------------------------------------
\subsection{Scaling de crat\`ere (Holsapple-Schmidt)}
% ----------------------------------------------------------------------------

\textbf{Source:} Holsapple, K. A. (1993).
\textit{The Scaling of Impact Processes in Planetary Sciences}.
Annu.\ Rev.\ Earth Planet.\ Sci., 21, 333--373.

Groupes $\pi$ adimensionnels:
\begin{align}
\pi_2 &= \frac{g \cdot a}{U^2} & \text{(gravit\'e vs inertie)} \\
\pi_3 &= \frac{Y}{\rho_t U^2} & \text{(r\'esistance vs inertie)} \\
\pi_4 &= \frac{\rho_t}{\rho_p} & \text{(densit\'es cible/projectile)} \\
\pi_V &= \frac{\rho_t V}{m} & \text{(volume crat\`ere adimensionnel)}
\end{align}

Loi de scaling compl\`ete:
\begin{equation}
\pi_V = K_1 \left[
\pi_2 \cdot \pi_4^{(6\nu - 2 - \mu)/(3\mu)}
+ K_2 \cdot \pi_3 \cdot \pi_4^{(6\nu - 2)/(3\mu)}
\right]^{-3\mu/(2+\mu)}
\end{equation}

R\'egime gravit\'e ($\pi_2 \gg \pi_3$, gros impacts):
\begin{equation}
\pi_V = K_1 \cdot \pi_2^{-3\mu/(2+\mu)} \cdot \pi_4^{(6\nu-2-\mu)/(2+\mu)}
\end{equation}

R\'egime r\'esistance ($\pi_3 \gg \pi_2$, petits impacts):
\begin{equation}
\pi_V = K_1 K_2^{-3\mu/(2+\mu)} \cdot \pi_3^{-3\mu/(2+\mu)} \cdot \pi_4^{(6\nu-2)/(2+\mu)}
\end{equation}

Exposant $\mu$: 2/3 (scaling \'energie) \`a 1/3 (scaling quantit\'e de mouvement).
Valeur exp\'erimentale pour mat\'eriaux consolid\'es: $\mu \approx 0.58$.

% ----------------------------------------------------------------------------
\subsection{Mapping Sedov-Taylor $\to$ Mycelium Yggdrasil}
% ----------------------------------------------------------------------------

\begin{center}
\begin{tabular}{lll}
\toprule
\textbf{Physique} & \textbf{Yggdrasil} & \textbf{Mesure} \\
\midrule
$E$ (\'energie) & strate\_height $\times$ continents\_touch\'es & entier \\
$\rho_0$ (densit\'e) & meshedness $\alpha$ locale avant impact & $[0, 1]$ \\
$R(t)$ (rayon choc) & concepts affect\'es \`a $t$ mois & entier \\
$\dot{R}(t)$ (vitesse) & nouvelles co-occurrences / mois & taux \\
$\gamma$ (adiabatique) & rigidit\'e du sol local & \`a calibrer \\
$\beta$ (constante) & \`a calibrer depuis bo\^ites de m\'et\'eorites & $\sim 1$ \\
\bottomrule
\end{tabular}
\end{center}

\textbf{Formule adapt\'ee:}
\begin{equation}
\boxed{R(t) = \beta_{\text{myc}} \left(\frac{E_{\text{impact}} \cdot t^2}{\alpha_{\text{local}}}\right)^{1/5}}
\end{equation}
avec:
\begin{align}
E_{\text{impact}} &= \text{strate\_height} \times \text{continents\_touch\'es} \\
\alpha_{\text{local}} &= \text{meshedness avant impact (densit\'e du sol S-2 \`a S0)}
\end{align}

\textbf{Pr\'ediction testable (corr\'elation bougie $\leftrightarrow$ type de trou):}
\begin{itemize}
\item Trou A (technique): $\alpha_{\text{local}}$ \'elev\'e $\to$ blast lent $\to$ bougie moyenne
\item Trou B (conceptuel): $\alpha_{\text{local}}$ faible $\to$ blast rapide $\to$ bougie courte
\item Trou C (perceptuel): $\alpha_{\text{local}}$ \'elev\'e + hostile $\to$ blast bloqu\'e $\to$ bougie longue
\end{itemize}

\textbf{Calibration:} depuis Shannon 1948 (premi\`ere bo\^ite, $\rho_0$ mesurable).
Test final: G\"odel 1931 (avant lui tout est plat, $\rho_0 \approx 0$).

% ============================================================================
\section{Bougies OHLC et 7 deltas}
% ============================================================================

\textbf{Fichier:} \texttt{engine/meteorites.py} (V3, cod\'e session 7, 24 f\'ev 2026)

Chaque perc\'ee majeure = une bougie OHLC:
\begin{center}
\begin{tabular}{ll}
\toprule
\textbf{OHLC} & \textbf{D\'efinition} \\
\midrule
Open & Date d'\'emission du paper \\
High & Pic de reconfiguration maximale du mycelium \\
Low & Creux (r\'esistance paradigme / stabilisation) \\
Close & Date de validation (prouv\'e, r\'epliqu\'e, explosion citations) \\
\bottomrule
\end{tabular}
\end{center}

Longueur de bougie $= t_{\text{Close}} - t_{\text{Open}}$ = temps de r\'esistance du paradigme.

\textbf{7 indicateurs techniques} (deltas entre frame avant et apr\`es):
\begin{align}
\Delta_1 &= \text{volume} = \text{nouvelles ar\^etes co-occurrence cr\'e\'ees} \\
\Delta_2 &= \text{amplitude} = \|\mathbf{x}_{\text{apr\`es}} - \mathbf{x}_{\text{avant}}\|_2
  \quad\text{(d\'eplacement spectral des centro\"\i des)} \\
\Delta_3 &= \Delta BC = BC_{\text{apr\`es}} - BC_{\text{avant}}
  \quad\text{(nouveaux ponts vs obsol\`etes)} \\
\Delta_4 &= \Delta\alpha = \alpha_{\text{apr\`es}} - \alpha_{\text{avant}}
  \quad\text{(r\'esilience r\'eseau)} \\
\Delta_5 &= \Delta P4 = P4_{\text{ferm\'es}} - P4_{\text{ouverts}}
  \quad\text{(trous net)} \\
\Delta_6 &= \Delta\text{Physarum} = \text{hyphes cr\'e\'ees} - \text{hyphes mortes} \\
\Delta_7 &= \text{births} = \text{concepts n\'es} - \text{concepts morts}
\end{align}

\textbf{Impact pond\'er\'e:}
\begin{equation}
I_{\text{m\'et\'eorite}} = \text{strate\_height} \times \text{continents\_touch\'es}
\end{equation}
Exemples: Poincar\'e ($S3 \times 1 = 3$), Shannon ($S1 \times 7 = 7$),
G\"odel ($S6 \times 9 = 54$, hors \'echelle).

% ============================================================================
\section{Co-occurrence et force de lien}
% ============================================================================

\textbf{Fichier:} \texttt{engine/core/scisci.py:128}

\begin{equation}
\text{CoOcc}(A,B) = \frac{P(A \cap B)}{P(A) \cdot P(B) \cdot N}
= \frac{n_{AB}}{(n_A / N) \cdot (n_B / N) \cdot N}
\end{equation}
$> 1$: sur-repr\'esent\'e. $< 1$: sous-repr\'esent\'e. $= 0$: jamais vu ensemble (trou P4).

\textbf{Jaccard modifi\'e:}
\begin{equation}
J(A,B) = \frac{n_{AB}}{n_A + n_B - n_{AB}}
\end{equation}

% ============================================================================
\section{Hypoth\`eses test\'ees (V3, sessions 33--35)}
% ============================================================================

\begin{enumerate}
\item \textbf{Partition \'energ\'etique:} NON TEST\'E. Sedov invalid\'e avant test.
\item \textbf{$\gamma$ variable:} NON APPLICABLE. Sedov invalid\'e.
\item \textbf{Exposant 2/5:} \textbf{INVALID\'E}. $R(t) \propto t^{2/5}$ ne tient pas.
  La logistique $R(t) = K/(1+e^{-r(t-t_0)})$ fitte mieux (R$^2$ = 1.00).
\item \textbf{Corr\'elation bougie $\leftrightarrow$ $\rho_0$:} \textbf{PARTIELLEMENT CONFIRM\'E}.
  Trou A $\to$ bougie longue, volume faible (Laser, $r=1.16$).
  Trou B $\to$ bougie courte, volume \'elev\'e (CRISPR, $r=5.41$).
  Trou C $\to$ bougie tr\`es longue (mRNA, 30 ans avant explosion).
\end{enumerate}

% ============================================================================
\section{Le caillou dans la mare (V3 refonte, session 33)}
% ============================================================================

Le myc\'elium $=$ une mare, la perc\'ee $=$ un caillou.
L'\'energie du caillou NE pr\'edit PAS $R_{\max}$. C'est la mare qui d\'ecide.

\textbf{\'Energie d'impact (INVALID\'E):}
\begin{equation}
E = m \times g \times h
\end{equation}
Test\'e avec $m = \text{works\_count}$, $h = \text{strate}$:
$\rho < 0.26$ pour toutes les formules. FAIL.

\textbf{R\'esultat cl\'e:} \textit{median\_neighbor\_works} (poids m\'edian des voisins
dans cooc\_global) pr\'edit $R_{\max}$ avec $\rho = +0.60$ ($p = 0.029$, $n = 13$).

\textbf{13 m\'et\'eorites mesur\'ees} (BFS sur WT3, 885M rows cooc):
mRNA 97\%, CRISPR 91\%, Higgs 85\%, Internet 82\%, AlphaFold 79\%,
Grav waves 77\%, Shannon 76\%, Transistor 70\%, Turing 65\%, ADN 55\%,
G\"odel 44\%, Poincar\'e 32\%, Laser 6\%.

% ============================================================================
\section{Logistique S-curve (trajectoire $R(t)$)}
% ============================================================================

\textbf{Source:} Verhulst, P. F. (1838).
\textit{Notice sur la loi que la population suit dans son accroissement}.
Corr.\ Math.\ et Phys., 10, 113--121.

\begin{equation}
\boxed{R(t) = \frac{K}{1 + e^{-r(t - t_0)}}}
\end{equation}
\begin{itemize}
\item $K$ = capacit\'e (port\'ee maximale $\approx R_{\max}$)
\item $r$ = taux de croissance (vitesse de l'onde)
\item $t_0$ = point d'inflexion
\end{itemize}

Test\'e sur 13 m\'et\'eorites: R$^2$ m\'edian $= 0.9996$, $K_{\text{error}} < 1\%$.
\textbf{PASS}.

% ============================================================================
\section{Mort spectrale (pr\'ediction de la dur\'ee de vie)}
% ============================================================================

\textbf{Source:} Chung, F. R. K. (1997).
\textit{Spectral Graph Theory}. AMS CBMS No.\ 92.

\begin{equation}
\boxed{t_{\text{death}} \approx \frac{1}{\lambda_1} \times \text{ratio}}
\end{equation}
\begin{itemize}
\item $\lambda_1 = 0.226$ = gap spectral du Laplacien WT4 (66,342 n\oe uds)
\item ratio $\approx 2.0$ yr/unit (fitt\'e sur 13 m\'et\'eorites)
\item Pr\'ediction: $t_{\text{death}} = 8.1$ ans
\item Observ\'e: 5--12 ans, MAE $= 1.79$ ans
\end{itemize}

Z\'ero param\`etre libre. Tout est dans la structure du graphe. \textbf{PASS}.

% ============================================================================
\section{Newman SIR-percolation (port\'ee $R_{\max}$)}
% ============================================================================

\textbf{Source:} Newman, M. E. J. (2002).
\textit{Spread of epidemic disease on networks}.
Phys.\ Rev.\ E, 66, 016128.
\href{https://doi.org/10.1103/PhysRevE.66.016128}{doi:10.1103/PhysRevE.66.016128}.
arXiv: \texttt{cond-mat/0205009}.

\begin{equation}
S(T) = 1 - G_0(u)
\end{equation}
o\`u $u$ r\'esout $u = G_1(1 - T + Tu)$:
\begin{align}
G_0(x) &= \sum_k p_k \cdot x^k & \text{(generating function du degr\'e)} \\
G_1(x) &= \frac{G_0'(x)}{G_0'(1)} & \text{(excess degree)} \\
T_c &= \frac{\langle k \rangle}{\langle k^2 \rangle - \langle k \rangle}
  & \text{(seuil de percolation)}
\end{align}

Sur Yggdrasil: $\langle k \rangle = 2136$, $T_c = 0.0002$.
Test temporel $K$ pr\'edit \`a $\pm 16\%$ (PASS).

% ============================================================================
\section{ETAS / Loi d'Omori (Carmack move: sismologie $\times$ \'epid\'emie)}
% ============================================================================

\textbf{Carmack move:} seismology $\times$ epidemic, cooc $= 0$, score $= 321.6$.
Ce mod\`ele n'avait JAMAIS \'et\'e appliqu\'e aux graphes de connaissances.

\textbf{Sources:}
\begin{itemize}
\item Omori, F. (1894). \textit{On the aftershocks of earthquakes}.
  J.\ College of Science, 7, 111--200.
\item Utsu, T. (1961). \textit{A statistical study on the occurrence of aftershocks}.
  Geophys.\ Mag., 30, 521--605.
\item Ogata, Y. (1988). \textit{Statistical models for earthquake occurrences}.
  JASA, 83, 9--27. \href{https://doi.org/10.1080/01621459.1988.10478560}{doi:10.1080/01621459.1988.10478560}
\item Saichev, A. \& Sornette, D. (2005). Phys.\ Rev.\ E, 71, 016608.
  \href{https://doi.org/10.1103/PhysRevE.71.016608}{doi:10.1103/PhysRevE.71.016608}
\end{itemize}

\textbf{Omori-Utsu (taux de r\'epliques):}
\begin{equation}
\boxed{R_{\text{rate}}(t) = \frac{K}{(t + c)^p}}
\end{equation}
Test\'e sur 13 m\'et\'eorites (new\_concepts/year): R$^2$ m\'edian $= 0.87$. \textbf{PASS}.

\textbf{ETAS branching sur $\mu(t)$:}
\begin{equation}
\mu(t) = \frac{A}{(t + c)^p} + \mu_{\text{bg}}
\end{equation}
R$^2$ m\'edian $= 0.998$. \textbf{PASS}.

\textbf{R\'esultat cl\'e: $p$ est UNIVERSEL.}
\begin{equation}
p = 4.74 \pm 0.63 \quad (\text{CV} = 13\%, n = 12)
\end{equation}
La loi d'Omori des r\'epliques sismiques s'applique aux perc\'ees scientifiques.

% ============================================================================
\section{Candlestick model (Carmack move: finance $\times$ science)}
% ============================================================================

\textbf{Sources:}
\begin{itemize}
\item Homma, M. (1755). Riz de Dojima --- premi\`ere utilisation des candlesticks.
\item Nison, S. (1991). \textit{Japanese Candlestick Charting Techniques}. Prentice Hall.
\item Sky, session 7 (24 f\'ev 2026). Premi\`ere formulation OHLC scientifique.
\end{itemize}

\textbf{Bougie scientifique OHLC:}
\begin{center}
\begin{tabular}{lll}
\toprule
\textbf{Composante} & \textbf{D\'efinition} & \textbf{Mesure} \\
\midrule
Open & Date de publication & $t = 0$ \\
High & Pic de reconfiguration (peak\_edges) & $t \approx 2$ ans \\
Low & Creux de r\'esistance du paradigme & variable \\
Close & Mort de l'onde & death\_t ($\approx 8.1$ ans) \\
Longueur & Close $-$ Open & death\_t \\
Volume & peak\_edges & mesur\'e \`a $t+2$ \\
\bottomrule
\end{tabular}
\end{center}

\textbf{Formule cl\'e:}
\begin{equation}
\boxed{\text{candle\_ratio} = \frac{\text{peak\_edges}}{\text{death\_t}}}
\end{equation}
\begin{equation}
r = 0.761 \times \log(\text{candle\_ratio}) - 5.116
\end{equation}

\textbf{R\'esultat:} candle\_ratio vs $r$: $\rho = +0.9176$ ($p < 0.0001$).
C'est le plus fort signal pour la vitesse de l'onde.

\textbf{Test temporel (train pr\'e-1960 $\to$ predict post-1974):}
R$^2$ trajectoire passe de $-0.09$ (sans candlestick) \`a $+0.44$ (avec).
LOO: $r$ erreur m\'ediane $= 18\%$, $t_0$ erreur m\'ediane $= 16\%$.

\textbf{Mod\`ele en 2 temps:}
\begin{enumerate}
\item Mode 1 (pr\'e-impact, $t=0$): $K$ depuis la mare, death depuis le gap spectral.
\item Mode 2 (early warning, $t+2$): mesurer peak\_edges $\to$
  candle\_ratio $\to$ $r$ $\to$ $R(t)$ complet.
\end{enumerate}

\textbf{Type de bougie $\leftrightarrow$ type de trou:}
\begin{center}
\begin{tabular}{llll}
\toprule
\textbf{Type} & \textbf{Bougie} & \textbf{$r$} & \textbf{Exemple} \\
\midrule
A (technique) & Longue, volume faible & $r$ faible & Laser ($r = 1.16$) \\
B (conceptuel) & Courte, volume \'elev\'e & $r$ \'elev\'e & CRISPR ($r = 5.41$) \\
C (perceptuel) & Tr\`es longue, retard\'ee & $r$ variable & mRNA (30 ans) \\
\bottomrule
\end{tabular}
\end{center}

% ============================================================================
\section{Score Carmack (d\'etection de moves cross-domaine)}
% ============================================================================

\textbf{Fichier:} \texttt{engine/analysis/scan\_carmack\_moves.py}

\begin{equation}
\text{carmack\_score} = \text{desert\_ratio} \times \log(1 + \text{works}) \times |z_{\text{avg}}|
\end{equation}
\begin{itemize}
\item desert\_ratio $=$ fraction de paires avec cooc $= 0$
\item Tiers: S (nucl\'eaire, $\geq 0.8$), A (fort, $\geq 0.5$), B (moyen, $\geq 0.3$), C (exp\'erimental)
\end{itemize}

\textbf{D\'eserts confirm\'es (session 35):}
\begin{center}
\begin{tabular}{lrl}
\toprule
\textbf{Paire} & \textbf{Score} & \textbf{Status} \\
\midrule
seismology $\times$ epidemic & 321.6 & ETAS valid\'e \\
heat\_kernel $\times$ cascade & 116.3 & \`a explorer \\
cognitive\_psych $\times$ scale\_free & 75.1 & \`a explorer \\
logistic\_function $\times$ scale\_free & 0 & notre mod\`ele \\
carrying\_capacity $\times$ scale\_free & 0 & \`a explorer \\
\bottomrule
\end{tabular}
\end{center}

% ============================================================================
\section{R\'esum\'e du mod\`ele V3 final}
% ============================================================================

\begin{center}
\begin{tabular}{lllll}
\toprule
\textbf{Composante} & \textbf{Formule} & \textbf{Params} & \textbf{Verdict} & \textbf{Source} \\
\midrule
Trajectoire $R(t)$ & $K/(1+e^{-r(t-t_0)})$ & 3 & PASS (R$^2$ = 1.00) & Verhulst 1838 \\
Mort & $1/\lambda_1 \times \text{ratio}$ & 0 & PASS (1.79 ans) & Chung 1997 \\
Port\'ee $K$ & Ridge(mare) & Ridge & PASS ($\pm 16\%$) & Newman 2002 \\
Vitesse $r$ & $0.76 \log(\text{ratio}) - 5.1$ & 0* & LOO 18\% & Candlestick \\
D\'ecroissance & $K/(t+c)^p$, $p = 4.74$ & 0 & PASS (R$^2$ = 0.87) & Omori 1894 \\
\'Energie $E$ & $m \times g \times h$ & 3 & FAIL & --- \\
Sedov-Taylor & $\beta(Et^2/\rho_0)^{1/5}$ & 2 & FAIL & Invalid\'e \\
\bottomrule
\end{tabular}
\end{center}

*$r$ d\'eriv\'e du candle\_ratio qui est mesur\'e \`a $t+2$, pas pr\'edit pr\'e-impact.

\textbf{Le mod\`ele:}
\begin{enumerate}
\item Un papier sort ($t = 0$). On mesure la mare (cooc\_global autour des seeds).
\item On pr\'edit $K \approx \pm 16\%$ et death $\approx 8.1$ ans.
\item Apr\`es 2 ans ($t + 2$), on mesure peak\_edges.
\item On calcule $r$ depuis candle\_ratio. On reconstruit $R(t)$.
\item La d\'ecroissance suit Omori avec $p = 4.74$ (universel).
\end{enumerate}

% ============================================================================
\section{Index centralis\'e --- Scripts, r\'esultats, donn\'ees}
% ============================================================================

\subsection{Scripts principaux}

\begin{center}
\begin{tabular}{lll}
\toprule
\textbf{Script} & \textbf{R\^ole} & \textbf{Session} \\
\midrule
\texttt{engine/core/scisci.py} & Fitness $\eta$, D-index, z-score, Q-factor & 6 \\
\texttt{engine/core/holes.py} & D\'etection trous A/B/C, score P4 & 6 \\
\texttt{engine/analysis/meteorites.py} & Sedov (invalid\'e), OHLC, 7 deltas, Frame & 7 \\
\texttt{engine/topology/species\_identifier.py} & 5 curseurs Lehmann, 9 esp\`eces & 8 \\
\texttt{engine/topology/species\_classifier.py} & Clustering spectral K=9 & 8 \\
\texttt{engine/topology/winter\_tree\_scanner.py} & WT1: concept$\times$concept, 108M paires & 10 \\
\texttt{engine/topology/wt2\_scanner.py} & WT2: per-paper, 832K papers & 22 \\
\texttt{engine/topology/wt3\_bible.py} & WT3: La Bible, 78 GB, 885M cooc & 23--27 \\
\texttt{engine/topology/wt4\_spectral.py} & WT4: Laplacien 66K n\oe uds, 75.6M ar\^etes & 28 \\
\texttt{engine/topology/impact\_scale.py} & Impact Scale 0--10 (833K papers) & 31 \\
\texttt{engine/analysis/scan\_carmack\_moves.py} & Score Carmack, d\'eserts cross-domaine & 30 \\
\texttt{engine/analysis/calibrate\_sedov.py} & Calibrage Sedov (invalid\'e session 33) & 32 \\
\texttt{scripts/godel\_holdout\_wave.py} & BFS onde 13 m\'et\'eorites & 33 \\
\texttt{scripts/wave\_comprehensive\_test.py} & Batterie 7 phases, 6 mod\`eles, 13 m\'et. & 34 \\
\texttt{scripts/wave\_predictive\_model.py} & R\'egression K, r, t$_0$ depuis mare & 34 \\
\texttt{scripts/wave\_predictive\_search.py} & Nouvelles features + Carmack search & 34 \\
\texttt{scripts/wave\_etas\_carmack.py} & Carmack ETAS/Omori sur 13 m\'et. & 35 \\
\texttt{scripts/wave\_candlestick\_model.py} & Mod\`ele candlestick, $\rho = 0.92$ & 35 \\
\texttt{scripts/wave\_hybrid\_final.py} & Mod\`ele hybride + test temporel & 35 \\
\texttt{scripts/scan\_openalex\_years.py} & Scan OpenAlex papers/an (693K) & 35 \\
\bottomrule
\end{tabular}
\end{center}

\subsection{R\'esultats JSON}

\begin{center}
\begin{tabular}{lll}
\toprule
\textbf{Fichier} & \textbf{Contenu} & \textbf{Taille} \\
\midrule
\texttt{wave\_comprehensive\_test.json} & 13 m\'et. $\times$ 7 phases $\times$ 6 mod\`eles & 43 KB \\
\texttt{wave\_test\_plan.json} & Batterie 5 tests initiale (6 m\'et.) & 9 KB \\
\texttt{wave\_model\_test.json} & P(k), Newman, eigenvalues, branching & 4 KB \\
\texttt{wave\_full\_test.json} & Propri\'et\'es seeds + mare (13 m\'et.) & 6 KB \\
\texttt{wave\_research\_v2.json} & 134 papers, 9 mod\`eles, P4 analysis & 17 KB \\
\texttt{wave\_predictive\_model.json} & R\'egression K, r, t$_0$ + G\"odel holdout & 17 KB \\
\texttt{wave\_predictive\_search.json} & Nouvelles features, Carmack d\'eserts & 12 KB \\
\texttt{wave\_etas\_carmack.json} & Omori/ETAS fits 13 m\'et., p universel & -- \\
\texttt{wave\_candlestick\_model.json} & Bougies OHLC, candle\_ratio vs $r$ & -- \\
\texttt{wave\_hybrid\_final.json} & Mod\`ele hybride, test temporel & -- \\
\texttt{scan\_carmack\_moves.json} & 34 moves, 23 techniques, d\'eserts & 60 KB \\
\texttt{scan\_wave\_propagation\_research.json} & 761 papers, 12 mod\`eles & 22 KB \\
\bottomrule
\end{tabular}
\end{center}

Tous dans \texttt{data/results/}.

\subsection{Donn\'ees sources}

\begin{center}
\begin{tabular}{lll}
\toprule
\textbf{Donn\'ee} & \textbf{Chemin} & \textbf{Taille} \\
\midrule
WT3 La Bible & \texttt{data/wt3.db} & 78 GB \\
WT4 Laplacien & \texttt{data/scan/wt4\_spectral.json} & 66K n\oe uds \\
65K concepts & \texttt{data/scan/concepts\_65k.json} & 7 MB \\
Positions spectrales & \texttt{data/scan/spectral\_births.json} & 66K pos. \\
arXiv sources & \texttt{E:/arxiv/src/} & 1 TB \\
OpenAlex & \texttt{E:/openalex/data/} & 692 GB \\
\bottomrule
\end{tabular}
\end{center}

\subsection{Documentation}

\begin{center}
\begin{tabular}{ll}
\toprule
\textbf{Document} & \textbf{R\^ole} \\
\midrule
\texttt{docs/reference/formulas.tex} & Ce fichier --- toutes les formules \\
\texttt{docs/reference/wave\_physics.md} & Formules onde + r\'esultats tests \\
\texttt{docs/SOL.md} & Source de v\'erit\'e inter-sessions \\
\texttt{docs/CHANGELOG.md} & Historique des sessions \\
\texttt{docs/TODO.md} & Prochaines \'etapes \\
\texttt{data/metaprompt\_wave\_research.md} & Prompt V2 pour deep research \\
\texttt{data/metaprompt\_predictive\_wave.md} & Prompt V3 pour pr\'ediction \\
\bottomrule
\end{tabular}
\end{center}

\end{document}
